% Copyright (c) 2026 Marcio Luiz Lima de Oliveira and V. el Hakim.
% Distribution and modification are subject to LPPL version 1.3c.
% SPDX-License-Identifier: LPPL-1.3c
% Part of the LPPL-covered work listed in MANIFEST.txt.
% See LICENSE and LICENSING.txt for the licence and scope.

\chapter{Symbols, equations and acronyms}\label{chap:04}
\renewcommand{\FigLocation}{Chapters/Chapter_4/fig/}
\renewcommand{\ChapLocation}{Chapters/Chapter_4/}
\begin{absbox}
Declare notation centrally and introduce it where the explanation reads naturally. The global list of symbols and the abbreviation list use first-use information across the whole thesis.
\end{absbox}

\section{Declare a command before using it}
Add a symbol to \texttt{DeclareSymbols.tex}, below the registry implementation.
Choose a descriptive command name containing letters; do not overwrite a
standard command such as \verb|\alpha| or \verb|\epsilon|.
\begin{lstlisting}
\DeclareThesisSymbol{sampleInput}{x}{Input value}
\DeclareThesisSymbol{sampleOutput}{y}{Output value}
\DeclareThesisSymbol{sampleSlope}{a}{Linear model coefficient}
\DeclareThesisSymbol{sampleOffset}{b}{Linear model offset}
\end{lstlisting}
These declarations already exist in the supplied file. Once registered,
\verb|\sampleInput| can be used in text or maths. A declaration in the preamble
is a programming step; the explanation of that variable in your prose may
come before or after its first equation. An undefined command still cannot be
used before it has been registered with TeX.

\section{First appearances inside an equation}
The next equation deliberately uses four registered symbols for the first time:
\begin{equation}
  \sampleOutput = \sampleSlope\sampleInput + \sampleOffset.
  \label{eq:linear}
\end{equation}
Here \sampleInput\ is the input, \sampleOutput\ is the output, \sampleSlope\ is
the slope and \sampleOffset\ is the offset. These explanations follow the
equation by choice. Each symbol links to the global list, and its list entry
links back to the first typeset occurrence recorded on a shipped page.
\begin{lstlisting}
\begin{equation}
  \sampleOutput = \sampleSlope\sampleInput + \sampleOffset.
  \label{eq:linear}
\end{equation}
Here \sampleInput\ is the input and \sampleOutput\ is the output.
\end{lstlisting}
Use an ordinary equation label alongside the symbol commands. Symbol tracking
does not require a \verb|\label| for each variable. The historical interface
\verb|\mathsymbols{4}{sampleInput}{x}{Input value}| also works when used
\emph{instead of} the corresponding new declaration. Its chapter argument is
ignored: the list is global and follows declaration order, not chapter order
or alphabetical order. Register each command once.

\section{Aligned equations, vectors and delimiters}
\begin{align}
  \sampleVector &= (x_1,\ldots,x_{\sampleCount})^{\mathsf T},\label{eq:vector}\\
  \sampleError &= \norm{\sampleMatrix\sampleVector-\sampleVector}{2}.
  \label{eq:residual}
\end{align}
The example in \eqref{eq:residual} assumes a compatible square matrix; its
purpose is to show typography. The dimension notation is explained here after
the display: \sampleCount\ is the number of observations, \sampleVector\ is
the observation vector, \sampleMatrix\ is a linear transformation and
\sampleError\ is the residual magnitude.
\begin{lstlisting}
\norm{x}{2}          % automatically sized norm, with subscript 2
\norm{x}{}           % norm without a subscript
\norm[\big]{x}{2}    % explicit delimiter size
\Abs{u-v}           % automatically sized absolute value
\Abs[\Big]{u-v}     % explicit size
\Set{x}{x>0}        % fixed-size set braces by default
\Set[\Big]{x}{x>0}  % larger set braces
\func{f}{X}{Y}      % f : X -> Y
\end{lstlisting}
An optional size is useful in multi-line displays because a \verb|\left|
delimiter cannot be opened on one alignment row and closed on another.
The legacy norm always takes two mandatory arguments, including an empty
second argument when no subscript is wanted. Do not write \verb|\norm{x}|.
\[
  \func{f}{\mathbb R}{\mathbb R},\qquad
  f(x)=\Abs{x},\qquad X=\Set{x\in\mathbb R}{x>0}.
\]
\texttt{MathOperators.tex} also provides upright operators such as
\verb|\argmin|, \verb|\argmax|, \verb|\diag| and \verb|\Tr|.
The notation appendix lists the full helper vocabulary. Use \verb|\mathrm{d}|
for a differential and a consistent vector convention throughout the thesis.

\section{What the list records}
Only used registered symbols appear in the list. The declared
\verb|\unusedExample| is deliberately never executed and should be absent.
Plain letters such as \verb|x| in an equation are ordinary maths, not tracked
symbols. A code listing does not count as a use. Table-of-contents replays,
glossary rows and temporary equation measurements do not create first-use
records. Float placement can affect which page is first, so let references
stabilise after changing figures or captions.

The list reflects the previous complete compilation. Rebuild the complete
thesis after deleting uses or working with \verb|\includeonly|. There is no
separate formal-definition command: the default keeps the existing first-use
behaviour. Avoid symbols in headings when plain words are clearer. If a heading
needs mathematics, supply readable bookmark text:
\begin{lstlisting}
\section{Behaviour of \texorpdfstring{$\sampleInput$}{the input}}
\end{lstlisting}

\section{Define and use acronyms}\label{sec:acronyms}
Define abbreviations in \texttt{Chapters/Chapter\_Appendix/Acronyms.tex}:
\begin{lstlisting}
\begin{acronym}[NASA]
  \acro{UT}{University of Twente}
  \acro{PDF}{Portable Document Format}
  \acro{API}{application programming interface}
\end{acronym}
\end{lstlisting}
The optional \texttt{NASA} specifies the width reserved for the short forms;
choose the widest abbreviation in your own list. Definitions located in the
back matter become available to earlier pages through auxiliary files, so an
initial run may show unresolved acronyms. Complete the usual rebuild sequence.

The first \verb|\ac{UT}| produces \ac{UT}; the next produces \ac{UT}.
For another example, \ac{API} becomes \ac{API} on later use. The main commands are:
\begin{description}
\item[\texttt{\string\ac\{key\}}] Full form on first use, short form afterwards.
\item[\texttt{\string\acs\{key\}}] Short form, such as \acs{PDF}.
\item[\texttt{\string\acl\{key\}}] Long form, such as \acl{PDF}.
\item[\texttt{\string\acf\{key\}}] Full form regardless of earlier use.
\item[\texttt{\string\acp\{key\}}] Plural, with the same first-use convention.
\item[\texttt{\string\ac*\{key\}}] An untracked use; a star is also available on
the explicit short, long and full forms.
\end{description}
For irregular plurals add \verb|\acroplural{key}[short plural]{long plural}|.
For example the sample \texttt{API} definition provides \acp{API}.
The options \texttt{printonlyused} and \texttt{withpage} select a global list of used
abbreviations with page information. \verb|\acreset{key}| or
\verb|\acresetall| resets first-use expansion if you intentionally want a fresh
introduction; the template does not reset the list at every chapter.
Remove \texttt{printonlyused} only if you want unused definitions printed too.

\section{Theorems, definitions and propositions}
\begin{definition}{Nonnegative input}{nonnegative}
An input is nonnegative when $x\geq 0$.
\end{definition}
\begin{theorem}{Nonnegative square}{square}
For every real number $x$, $x^2\geq 0$.
\end{theorem}
\begin{proof}
If $x\geq0$, the product $x\cdot x$ is nonnegative. If $x<0$, it is the
product of two negative numbers and is positive.
\end{proof}
\begin{proposition}{Equality case}{zero}
A real number satisfies $x^2=0$ if and only if $x=0$.
\end{proposition}
\begin{lstlisting}
\begin{theorem}{Nonnegative square}{square}
  For every real number $x$, $x^2\geq 0$.
\end{theorem}
Theorem~\ref{thrm:square}
\end{lstlisting}
These boxed environments take a title and a label suffix. Their automatic
prefixes are \texttt{thrm:}, \texttt{def:} and \texttt{prop:}; the suffix alone
is not a complete reference. They are numbered within sections. The
\texttt{definition} box can break across pages; keep the other two short or add
\texttt{breakable} to their \texttt{newtcbtheorem} styles when needed.
The boxed mathematical proposition is independent of the propositions insert.
